
\documentclass[11pt]{article}
\usepackage[margin=1in]{geometry}
\usepackage{amsmath,amssymb,booktabs,hyperref,graphicx,parskip,enumitem}
\usepackage{xcolor}
\hypersetup{colorlinks=true,linkcolor=blue!50!black,urlcolor=blue!50!black}
\setlist{nosep}

\title{Replication Report: (title not recorded)\\[2pt]\large OSTI 1868518 \textbar{} Coverage 5/10, Agreement 5/10}
\author{Rick Stevens \and Ollie (AI Assistant)\\\small Argonne National Laboratory}
\date{2026-04-27}

\begin{document}\maketitle

\section{Source paper}
\begin{itemize}
\item \textbf{Title:} (title not recorded)
\item \textbf{Authors:} Not recorded
\item \textbf{Year:} n/a
\item \textbf{Domain:} Not recorded
\item \textbf{Identifier:} OSTI 1868518
\end{itemize}


\section{Replication objective}
The central claim of the paper concerns not recorded. Our objective was to reproduce the headline numerical/qualitative results within the constraints of the AI-driven Replication Project (24--72\,h, commodity GPU/CPU compute, no author contact). 

\section{Methodology}
We followed the standard Replication Project workflow: ingest the source PDF, extract method and data pointers, draft a replication plan, attempt the smallest end-to-end pipeline that produces a comparable number, then scale up where compute and data permit. Tools used in this replication are catalogued in the meta-paper's computational tools inventory (see \texttt{COMPUTATIONAL\_TOOLS\_INVENTORY.md}).

This paper went through the Tier-Lift pipeline; supplementary notes at \texttt{.openclaw/workspace/24h-progress/tier-lift-v2/1868518.md}.

\subsection*{Paper-directory README excerpt}
\small
\par\medskip\textbf{Replication: Learning Sequential Distribution System Restoration via Graph-Reinforcement Learning}\par

**OSTI ID:** 1868518  
**Authors:** Tianqiao Zhao, Jianhui Wang  
**Published:** IEEE Transactions on Power Systems, vol. 37, no. 2, pp. 1601-1611, March 2022  
**DOI:** 10.1109/TPWRS.2021.3102345  
**Affiliations:** Brookhaven National Laboratory (BNL), Southern Methodist University

\par\medskip\textbf{\# Summary}\par

This paper proposes a multi-agent Graph-Reinforcement Learning framework for sequential distribution system restoration after extreme events. The approach combines Graph Convolutional Networks (GCNs) with Double DQN to coordinate distributed generators (DGs) and controllable switches. Each DG is modeled as an RL agent with parameter sharing for scalability. The GCN encodes the power network topology to capture device interactions.

\par\medskip\textbf{\# Replication}\par

\par\medskip\textbf{\#\# What We Replicated}\par
\par$\bullet$\ GCN + Double DQN (Dueling) multi-agent architecture
\par$\bullet$\ MLP-DQN baseline (no graph structure)  
\par$\bullet$\ Random and greedy heuristic baselines
\par$\bullet$\ IEEE 33-bus and IEEE 123-bus test networks
\par$\bullet$\ pandapower-based power flow environment
\par$\bullet$\ Training, evaluation, and comparison

\par\medskip\textbf{\#\# Key Findings}\par
1. **GCN-DQN learns effective restoration policies** — achieving up to 0.42 load restoration ratio on IEEE 33-bus (peak evaluation)
2. **GCN and MLP show similar performance** on small networks (33-bus) with 2,000 training episodes
3. **Larger networks are harder** — IEEE 123-bus achieves lower load ratios (\textasciitilde{}0.14 peak), consistent with increased coordination difficulty
4. **The framework is computationally accessible** — all training completes in \textless{}2 hours on CPU

\par\medskip\textbf{\#\# Files}\par
[code block omitted]

\par\medskip\textbf{\#\# Reproducibility}\par
\par$\bullet$\ **Software:** Python 3.12, PyTorch 2.2.2, pandapower 3.4.0, NetworkX 3.6.1
\par$\bullet$\ **Hardware:** Apple iMac, Intel x86\_64, CPU-only
\par$\bullet$\ **Training time:** \textasciitilde{}2 hours total (all methods, both networks)
\par$\bullet$\ **Seed:** 42 (reproducible)

\par\medskip\textbf{\# Status: ✅ Complete}\par
\normalsize

The replication was driven by an OpenClaw agent loop (Claude Opus / GPT-5 class models) issuing shell, Python, and (where applicable) GPU jobs. Compute usage and wall-time for individual papers are summarised in the meta-paper's resource table; not separately recorded for this paper unless noted in the Tier-Lift artefact. Sub-agents were spawned for replication-plan generation and (for papers in batches 1--4) for tier-lift extension runs.

\section{What was reproduced}
\begin{itemize}
\item Not recorded.
\end{itemize}

\section{What was missing or incomplete}
\begin{itemize}
\item Not recorded.
\end{itemize}
The dominant blocker categories for this paper, mapped onto the meta-paper's 9-category friction taxonomy (\S4), are listed in \S\ref{sec:friction} below.

\section{Quantitative agreement}
Per-quantity numerical comparisons are not separately tabulated in the structured evaluation record for this paper beyond the agreement rationale below; where the rationale references specific numbers, they are reproduced verbatim. A full numerical comparison table is recorded in the per-replication artefacts (see \S\ref{sec:repro}). For papers below 8/8, the agreement rationale identifies the specific quantities that diverge.

\paragraph{Agreement rationale.} Not recorded.

\section{Score and friction-point tagging}\label{sec:friction}
\begin{itemize}
\item \textbf{Coverage:} 5/10 --- Not recorded.
\item \textbf{Agreement:} 5/10 --- Not recorded.
\end{itemize}

\noindent\textbf{Friction points encountered (taxonomy tags):}
\begin{itemize}
\item \textbf{F5}: Force-field / hyperparameter drift (unspecified configuration)
\end{itemize}

\section{Follow-on questions}
\begin{itemize}
\item Not recorded.
\end{itemize}

\section{Reproducibility pointers}\label{sec:repro}
\begin{itemize}
\item \textbf{Paper directory:} \texttt{1868518-Graph-RL-Distribution-Restoration}
\item \textbf{Replication artefacts:}
\begin{itemize}
\item \texttt{/Users/stevens/Dropbox/REPLICATE-PROJECT/1868518-Graph-RL-Distribution-Restoration/replication}
\end{itemize}
\item \textbf{Replication-dir hint from JSONL:} \texttt{/Users/stevens/Dropbox/REPLICATE-PROJECT/1868518-Graph-RL-Distribution-Restoration}
\item \textbf{Status:} tier\_lift\_v2\_2026\_04\_27
\end{itemize}

To re-run, change directory into the paper directory above and consult its \texttt{README.md} for the entry-point script. Most replications follow the convention \texttt{replication/run\_*.sh} or \texttt{replication/<subdir>/run.py}.

\end{document}
